\begin{theorem}[最小外置预算的不可判定性]\label{thm:conclusion-minimal-external-budget-undecidable}
设 \(D,Y\) 为可数集合，并设 \(\pi:D\to Y\) 为全定义可计算映射，且具有有限纤维上界
\[
D(\pi):=\sup_{y\in Y}\#\pi^{-1}(y)<\infty.
\]
定义 \(\pi\) 的最小外置预算为
\[
b(\pi):=\min\Bigl\{B\in\NN_{\ge 1}:\ \exists\ \widetilde\pi:D\hookrightarrow Y\times\{0,1,\dots,B-1\}\ \text{单射},\ \operatorname{pr}_1\circ\widetilde\pi=\pi\Bigr\}.
\]
则：
\begin{enumerate}
  \item 对任意满足 \(D(\pi)<\infty\) 的 \(\pi\)，有
  \[
  b(\pi)=D(\pi).
  \]
  \item 存在一族显式可计算的投影算法 \(\{\pi_{M,x}\}\)（由图灵机 \(M\) 与输入 \(x\) 参数化）满足
  \[
  b(\pi_{M,x})=
  \begin{cases}
  1,& M(x)\ \text{不停机},\\
  2,& M(x)\ \text{停机},
  \end{cases}
  \]
  从而判定命题 \(b(\pi)=1\) 在该类算法上不可判定。
\end{enumerate}
\end{theorem}

\begin{proof}
第一条：记 \(d:=D(\pi)=\max_{y}\#\pi^{-1}(y)\)（有限性保证最大值存在）。
若 \(B<d\)，取 \(y\) 使 \(\#\pi^{-1}(y)>B\)，则任意满足 \(\operatorname{pr}_1\circ\widetilde\pi=\pi\) 的 \(\widetilde\pi\) 在纤维 \(\pi^{-1}(y)\) 上第二坐标仅有 \(B\) 个取值，鸽巢原理给出 \(\widetilde\pi\) 不可能单射，故 \(b(\pi)\ge d\)。
反之取 \(B=d\)，对每个纤维 \(\pi^{-1}(y)\) 固定一个注入编号 \(\mathrm{idx}:\pi^{-1}(y)\hookrightarrow\{0,\dots,d-1\}\)，并定义
\(
\widetilde\pi(x):=(\pi(x),\mathrm{idx}(x))
\)，即得单射提升，故 \(b(\pi)\le d\)。于是 \(b(\pi)=d\)。
\par
第二条：对给定 \(M,x\)，令 \(H_{M,x}(n)\) 表示谓词“图灵机 \(M\) 在输入 \(x\) 上于 \(n\) 步内停机”。
对每个 \(n\) 该谓词可计算（模拟 \(n\) 步即可）。
取 \(D:=\NN\times\{0,1\}\)，并定义全可计算映射 \(\pi_{M,x}:D\to\NN\)：
\[
\pi_{M,x}(n,b):=
\begin{cases}
4n+2b,& H_{M,x}(n)=0,\\
2n+1,& H_{M,x}(n)=1.
\end{cases}
\]
若 \(M(x)\) 不停机，则 \(H_{M,x}(n)=0\) 对所有 \(n\) 成立，故 \(\pi_{M,x}(n,b)=4n+2b\) 为单射，从而 \(b(\pi_{M,x})=1\)。
若 \(M(x)\) 于时刻 \(T\) 停机，则当 \(n\ge T\) 时 \(H_{M,x}(n)=1\)，于是 \(\pi_{M,x}(n,0)=\pi_{M,x}(n,1)=2n+1\)，故最大纤维大小至少为 \(2\)。
另一方面，由输出奇偶性分离可知跨分支不发生碰撞，且停机分支中每个奇数 \(2n+1\)（\(n\ge T\)）恰有两个原像 \((n,0),(n,1)\)，故最大纤维大小恰为 \(2\)，即 \(b(\pi_{M,x})=2\)。
因此 \(b(\pi_{M,x})=1\) 当且仅当 \(M(x)\) 不停机，故该判定问题不可判定。证毕。
\end{proof}

\begin{corollary}[预算函数的不可计算性与非半可判定性]\label{cor:conclusion-external-budget-noncomputable-nonre}
在上述有限纤维全定义可计算投影类上，映射 \(\pi\mapsto b(\pi)\) 不可计算；并且集合
\[
\{\pi:\ b(\pi)=1\}
\]
不是递归可枚举集。
\end{corollary}

\begin{proof}
若 \(b(\cdot)\) 可计算，则可判定 \(b(\pi)=1\)，与定理 \ref{thm:conclusion-minimal-external-budget-undecidable} 矛盾。
若集合 \(\{\pi:\ b(\pi)=1\}\) 为递归可枚举，则由同一定理中的构造可递归可枚举集合 \(\{(M,x):M(x)\text{不停机}\}\)，与经典结论矛盾。证毕。
\end{proof}

\endinput
